%!TEX root = ./main.tex

\section[CSI]{CSI: The Measurement You Already Have}

% SECTION TIMING: 8:00--22:00 (4 frames).
% Job of the section: put ONE object on the table -- the matrix H[antenna,
% subcarrier] -- and ONE enemy -- wrapping. Everything after is reading that
% matrix along one axis or the other. Students met subcarriers and training
% fields in Class 8; lean on that, do not re-teach OFDM.

% Teaching note (210 s): start from decoding, not from sensing. Ask: the wave
% arrives scaled and rotated -- how does the receiver know what "1+0j" was supposed
% to look like? Answer: the preamble. H falls out as a by-product on every packet.
\begin{frame}{The receiver must learn H before bit one}
  \vspace{0.3em}
  {\small Every packet opens with training symbols the receiver already knows
   (Class 8).\par}
  \vspace{0.3em}
  {\small On each subcarrier $f$ and each antenna, the receiver solves\par}
  \vspace{0.2em}
  \centering
  {\large $y = H \cdot x_{\text{known}} \;\;\Rightarrow\;\;
    H = y / x_{\text{known}}$\par}
  \vspace{0.6em}
  \raggedright
  {\small $H$ is one \textbf{complex number}: $|H|$ says how much the wave faded,
   $\arg H$ says\par}
  {\small how far it travelled --- for a single path of length $r$:\par}
  \vspace{0.2em}
  \centering
  {\large $H = \alpha\, e^{-j 2\pi f\, r/c}$
   \quad{\small\textcolor{softgray}{(one turn of phase per wavelength of $r$)}}\par}

  \vspace{0.7em}
  \qa{What does the modem do with $H$ after equalizing the packet?}
     {Throws it away --- thousands of times per second. CSI tools simply
      \emph{keep} it. Sensing is recycling.}

  \glossline{CSI = Channel State Information: the collection of these $H$ values
    \quad $c$ = speed of light}
\end{frame}

% Teaching note (180 s): the star figure of the lecture -- it returns in the SpotFi
% section with arrows on it. Point at one cell: one complex number. Then sweep a row,
% then a column. Do NOT yet say what the twists mean; just make them see the grid.
\begin{frame}{Ninety numbers per packet, for free}
  \vspace{0.3em}
  \centering
  \begin{tikzpicture}[scale=0.98,transform shape]
    \csigrid
    \node[font=\scriptsize,align=left,text=softgray] at (7.9,1.1)
      {each cell: one\\complex $H[m,k]$\\(arrow = its phase)};
  \end{tikzpicture}

  \vspace{0.5em}
  \raggedright
  {\small
   \begin{itemize}\itemsep3pt
     \item The classic Intel 5300 tool exports $3$ antennas $\times$ $30$
       subcarriers $=$ \textbf{90 complex values per packet}.
     \item Hundreds of packets per second $\Rightarrow$ a \emph{movie} of the
       channel: $H[\text{time},\, \text{antenna},\, \text{subcarrier}]$.
   \end{itemize}\par}

  {\scriptsize\color{softgray}Phase twists exaggerated on the frequency axis so
   they are visible.\par}

  \source{https://dhalperi.github.io/linux-80211n-csitool/}{Halperin et al., Linux 802.11n CSI Tool \quad$\cdot$\quad Tsinghua Wireless Sensing Tutorial, ``Understanding CSI''}
\end{frame}

% Teaching note (210 s): introduce the running example here and never change it:
% direct path 5 m at +30 deg (green), wall echo 8 m at -40 deg (purple). The formula
% is shown ONCE, then immediately translated into the three per-path words.
\begin{frame}{The room writes itself into the matrix}
  \vspace{0.2em}
  \centering
  \begin{tikzpicture}[scale=0.9,transform shape]
    % the room
    \draw[softgray!60] (0,2.6) -- (8.2,2.6);
    \node[font=\tiny,text=softgray] at (7.6,2.8) {wall};
    % AP with 3 antennas
    \node[draw=ranblue,fill=blue!5,rounded corners=2pt,minimum width=1.25cm,
          minimum height=0.6cm,font=\bfseries\small] (ap) at (0.7,0.3) {AP};
    \anttri{(0.35,0.75)}{ranblue}
    \anttri{(0.70,0.75)}{ranblue}
    \anttri{(1.05,0.75)}{ranblue}
    % phone
    \node[draw=softgray,fill=gray!10,rounded corners=1.5pt,minimum width=0.5cm,
          minimum height=0.85cm,font=\scriptsize] (ph) at (6.7,1.45) {\phantom{.}};
    \node[font=\scriptsize,text=softgray] at (6.7,2.0) {phone};
    % direct path
    \draw[-{Latex[length=2.2mm]},very thick,softgreen] (ph.west) -- (1.15,0.78);
    \node[font=\scriptsize\bfseries,text=softgreen] at (3.6,0.72)
      {direct: 5 m, from $+30^\circ$};
    % wall echo
    \draw[-{Latex[length=2.2mm]},thick,ranpurple] (ph.north) -- (3.9,2.6)
      -- (1.0,0.9);
    \node[font=\scriptsize\bfseries,text=ranpurple] at (5.6,2.33)
      {echo: 8 m, from $-40^\circ$};
  \end{tikzpicture}

  \vspace{0.4em}
  {\small Each path $\ell$ stamps the whole matrix at once:\par}
  \vspace{0.15em}
  \centering
  {\normalsize $H[m,k] \;=\; \sum_{\ell}\; \alpha_\ell\;
    \underbrace{e^{-j 2\pi f_k \tau_\ell}}_{\textcolor{coreorange}{\text{delay
    } \tau_\ell}}\;
    \underbrace{e^{-j \frac{2\pi}{\lambda} m d \sin\theta_\ell}}_{%
    \textcolor{ranblue}{\text{angle } \theta_\ell}}$\par}

  \vspace{0.45em}
  \takeaway{The AP receives only the \textbf{sum}. Sensing is un-summing:\\
    recover each path's $(\theta_\ell, \tau_\ell)$ from the matrix.}
\end{frame}

% Teaching note (180 s): the enemy, named once, properly. The odometer picture does
% the work: you see the dial, not the mileage. Close with the roadmap takeaway and
% point at the two colored axes -- that IS the table of contents of the lecture.
\begin{frame}{The catch: every phase arrives wrapped}
  \vspace{0.3em}
  \begin{columns}[T,onlytextwidth]
    \column{0.36\textwidth}
      \centering
      \begin{tikzpicture}[scale=0.85,transform shape]
        \draw[softgray] (0,0) circle (1.05);
        \draw[softgray!50] (-1.2,0) -- (1.2,0);
        \draw[softgray!50] (0,-1.2) -- (0,1.2);
        \draw[-{Latex[length=2.4mm]},very thick,ranblue] (0,0) -- (50:1.05);
        \node[font=\small\bfseries,text=ranblue] at (50:1.45) {$50^\circ$};
        \node[font=\scriptsize,text=softgray,align=center] at (0,-1.65)
          {what you measure:\\the needle, not the turns};
      \end{tikzpicture}
    \column{0.60\textwidth}
      {\small Our 5 m direct path at $f = 5.2$ GHz has turned\par}
      \vspace{0.15em}
      {\small \centering $f \cdot \tau \;\approx\; 5.2\,\text{GHz} \times 17\,
        \text{ns} \;\approx\; \text{\textbf{87 full turns}}$.\par}
      \vspace{0.4em}
      {\small The receiver sees only the leftover fraction ---
       like an odometer showing nothing but its last digit.\par}
      \vspace{0.4em}
      {\small One phase reading says: ``the path is $x + n\cdot 6$ cm,
       for \emph{some} integer $n$.'' Useless alone.\par}
  \end{columns}

  \vspace{0.6em}
  \takeaway{The fix is never one reading --- it is a \textbf{difference} along
    an axis:\\across \textcolor{ranblue}{antennas} $\to$ \textbf{angle};\;
    across \textcolor{coreorange}{subcarriers} $\to$ \textbf{distance}.}
\end{frame}
